% Compilar desde la raíz del repositorio.
\documentclass[main.tex]{subfiles}

\begin{document}
\graphicspath{{imagenes/tema1/}{../imagenes/tema1/}}

% Diapositiva original 1
\portadatema{1}{Funciones de varias variables}
% Diapositiva original 2
\begin{frame}{Funciones de varias variables}
La temperatura $T$ en un punto de la superficie de la Tierra a una hora determinada depende de la longitud y la latitud de dicho punto.
\medskip

Se puede pensar en $T$ como una función de dos variables y, si las denominamos $x$ e $y$, esta dependencia funcional se puede poner como
\[T=f(x,y).\]
\end{frame}
% Diapositiva original 3
\begin{frame}{Funciones de varias variables}
El volumen de una caja rectangular de dimensiones $x,y,z$ es $\text{Volumen}=x\cdot y\cdot z$.
\medskip

Este es un ejemplo de una función real de tres variables reales, que simbolizamos por:
\[f(x,y,z)=x\cdot y\cdot z.\]
\figura{s3_2.png}{.45}{.43}
\end{frame}
% Diapositiva original 4
\begin{frame}{Funciones de varias variables}
\textbf{Definición}
\medskip

En general, una función real de $n$ variables reales es una correspondencia que asigna a cada $n$-tupla $(x_1,x_2,\ldots,x_n)$ de su dominio un único valor
\[u=f(x_1,x_2,\ldots,x_n).\]
\end{frame}
% Diapositiva original 5
\begin{frame}{Funciones de varias variables}
Los conceptos conocidos para funciones de una variable tienen su equivalente para funciones de $n$ variables.
\end{frame}
% Diapositiva original 6
\begin{frame}{Funciones de varias variables}
Particularizando a dos variables, una función $f$ de dos variables asigna a cada par ordenado de números reales $(x,y)$ perteneciente a un conjunto $D$ un número real único denotado por $u=f(x,y)$.
\medskip

El conjunto $D$ es el dominio de $f$ y su imagen es el conjunto de valores que toma $f$.
\figura{s6_0.png}{.90}{.30}
\end{frame}
% Diapositiva original 7
\begin{frame}{Funciones de varias variables}
Las funciones de dos variables se suelen escribir de la forma $z=f(x,y)$ para explicitar el valor tomado por la función $f$ en un punto $(x,y)$.
\medskip

Las variables $x$ e $y$ son las variables independientes y $z$ es la variable dependiente.
\end{frame}
% Diapositiva original 8
\begin{frame}{Funciones de varias variables}
\small Si $f$ es una función de dos variables con dominio $D$, entonces la gráfica de $f$ es el conjunto de ternas $(x,y,z)$ de $\mathbb{R}^3$ tales que $z=f(x,y)$ y $(x,y)\in D$.
\medskip

Igual que las gráficas de funciones de una variable son curvas en el plano, las gráficas de funciones de dos variables son superficies en el espacio.
\figura{s8_2.png}{.44}{.47}
\end{frame}
% Diapositiva original 9
\begin{frame}{Funciones de varias variables}
Por ejemplo, $z=f(x,y)=x+y^2$ es una función de dos variables con dominio todos los pares de números reales y cuya representación en 3D es
\figura{s9_2.png}{.85}{.58}
\end{frame}
% Diapositiva original 10
\begin{frame}{Funciones de varias variables}
\textbf{Ejemplo}
\[z=f(x,y)=\frac{3x-5y}{y-x^2},\quad D=\{(x,y)\in\mathbb{R}^2:y-x^2\ne0\}.\]
\figura{s10_4.png}{.70}{.44}
El dominio es todo el plano $\mathbb{R}^2$ excepto los puntos de la parábola $y=x^2$.
\end{frame}
% Diapositiva original 11
\begin{frame}{Funciones de varias variables}
\textbf{Ejemplo}
\[z=f(x,y)=\sqrt{9-x^2-y^2},\quad D=\{(x,y)\in\mathbb{R}^2:9-x^2-y^2\ge0\}.\]
\figura{s11_5.png}{.77}{.40}
\small Como $x^2+y^2=9$ es la ecuación de una circunferencia de centro el origen y radio $3$, $D$ representa el interior y la frontera de dicha circunferencia.
\end{frame}
% Diapositiva original 12
\begin{frame}{Funciones de varias variables}
El dominio de la función
\[f(x,y,z)=\ln(z-y)+xy\sin z\]
es
\[D=\{(x,y,z)\in\mathbb{R}^3:z>y\}.\]
Es decir, todos los puntos que están por encima del plano $z=y$.
\end{frame}
% Diapositiva original 13
\begin{frame}{Funciones de varias variables}
Las funciones de varias variables pueden operarse de igual forma que las de una variable, ya sea con suma, producto, cociente o composición de funciones.
\end{frame}
% Diapositiva original 14
\begin{frame}{Curvas de nivel}
Las curvas de nivel de una función $f$ de dos variables son las curvas con ecuaciones $f(x,y)=k$, donde $k$ es una constante que pertenece a la imagen de $f$.
\begin{columns}[c]
\begin{column}{.48\textwidth}\figura{s14_3.png}{.85}{.56}\end{column}
\begin{column}{.48\textwidth}\figura{s14_5.png}{.98}{.56}\end{column}
\end{columns}
\end{frame}
% Diapositiva original 15
\begin{frame}{Curvas de nivel}
\textbf{Ejemplo}
\vfill
Obtener las curvas de nivel de
\[z=\sqrt{100-x^2-y^2}.\]
\vfill
\end{frame}
% Diapositiva original 16
\begin{frame}{Curvas de nivel}
\small\textbf{Ejemplo}\quad $z=\sqrt{100-x^2-y^2}$.
\medskip

Sustituimos la variable $z$ por una constante $z=c$:
\[c=\sqrt{100-x^2-y^2}.\]
Elevando al cuadrado,
\[c^2=100-x^2-y^2.\]
Reorganizando,
\[x^2+y^2=100-c^2.\]
Se obtienen circunferencias de centro el origen para $0\le c<10$. Para $c=10$ se obtiene el punto $(0,0)$. Para $c<0$ o $c>10$ no hay curvas de nivel.
\end{frame}
% Diapositiva original 17
\begin{frame}{Curvas de nivel}
\textbf{Ejercicio}\par
Obtener las curvas de nivel de la función de dos variables
\begin{columns}[c]
\begin{column}{.43\textwidth}
\[z=\sqrt{100-x^2-y^2}\]
\[x^2+y^2=100-c^2\]
\end{column}
\begin{column}{.55\textwidth}\figura{s17_3.png}{1}{.58}\end{column}
\end{columns}
\end{frame}
% Diapositiva original 18
\begin{frame}{Curvas de nivel}
\figura{s18_0.jpg}{.78}{.83}
\end{frame}
% Diapositiva original 19
\begin{frame}{Curvas de nivel}
\figura{s19_1.jpg}{1}{.72}\begin{center}\scriptsize Figure 19\end{center}
\end{frame}
% Diapositiva original 20
\begin{frame}{Level Curves}
\figura{s20_1.jpg}{1}{.76}
\end{frame}
% Diapositiva original 21
\begin{frame}{Curvas de nivel}
\figura{s21_0.jpg}{.73}{.86}
\end{frame}
% Diapositiva original 22
\begin{frame}{Funciones de varias variables}
\small\textbf{Ejercicios}
\medskip

Encontrar el dominio de la función
\[f(x,y,z)=\arcsin(xy)\,e^{3z}+e^{(y+z)\ln(x+z)}.\]
Describir los conjuntos de nivel para los valores de $k$ indicados:
\begin{enumerate}[a)]
\item $f(x,y)=x^2+y^2$,\quad $k=0,1,2,3$.
\item $f(x,y,z)=x^2+y^2$,\quad $k=0,1,2$.
\item $f(x,y)=x^2-y^2$,\quad $k=-2,-1,0,1,2$.
\end{enumerate}
En a) y c) se trata de curvas de nivel; en b), de superficies de nivel.
\end{frame}
\end{document}
